\section{Template Python 3}

Usar Python cuando la complejidad tiene margen, predominan enteros grandes,
strings, diccionarios o implementaci\'on corta. En problemas ajustados de
$O(n\log n)$ pesado, grafos enormes o DP con millones de estados, preferir C++.

\begin{lstlisting}[style=py]
import sys
from collections import Counter, defaultdict, deque
from functools import lru_cache
from heapq import heappop, heappush
from bisect import bisect_left, bisect_right
from math import gcd, isqrt, lcm, inf

input = sys.stdin.buffer.readline

def ints():
    return map(int, input().split())

def solve():
    n = int(input())
    values = list(ints())
    assert len(values) == n
    print(sum(values))

if __name__ == "__main__":
    solve()
\end{lstlisting}

Atajos importantes:
\begin{itemize}
  \item \texttt{math.isqrt(n)} da $\lfloor\sqrt n\rfloor$ exactamente.
  \item \texttt{pow(a,b,mod)} hace potencia modular y
  \texttt{pow(a,-1,mod)} calcula el inverso si existe.
  \item \texttt{sys.setrecursionlimit(...)} s\'olo ampl\'ia el l\'imite l\'ogico;
  para DFS muy profundo suele ser m\'as seguro usar una pila iterativa.
  \item Evitar \texttt{pop(0)} en listas; usar \texttt{deque.popleft()}.
  \item Usar input().decode() si se toma como b antes del input
\end{itemize}
